\documentclass[11pt,a4paper]{article}
\usepackage[utf8]{inputenc}
\usepackage[margin=1in]{geometry}
\usepackage{xcolor}
\usepackage{listings}
\usepackage{tcolorbox}
\tcbuselibrary{skins,breakable}
\usepackage{hyperref}
\usepackage{fancyhdr}
\usepackage{amsmath}

% Define premium colors
\definecolor{primary}{RGB}{31, 78, 121}    % Deep Blue
\definecolor{secondary}{RGB}{90, 107, 124} % Slate Gray
\definecolor{accent}{RGB}{235, 104, 65}    % Warm Coral/Orange
\definecolor{bglight}{RGB}{245, 247, 250}   % Off-white/Light blue-gray
\definecolor{codebg}{RGB}{240, 242, 245}    % Code block background
\definecolor{border}{RGB}{210, 215, 223}

% Document settings
\hypersetup{
    colorlinks=true,
    linkcolor=primary,
    filecolor=accent,
    urlcolor=primary,
}

% Code listing settings
\lstset{
    language=C++,
    backgroundcolor=\color{codebg},
    basicstyle=\ttfamily\small\color{black},
    keywordstyle=\color{primary}\bfseries,
    stringstyle=\color{accent},
    commentstyle=\color{secondary}\itshape,
    numbers=left,
    numberstyle=\tiny\color{secondary},
    stepnumber=1,
    numbersep=8pt,
    showstringspaces=false,
    breaklines=true,
    frame=single,
    rulecolor=\color{border},
    tabsize=4
}

% Custom boxes
\newtcolorbox{explanationbox}[1]{
    colback=bglight,
    colframe=primary,
    fonttitle=\bfseries\color{white},
    title=#1,
    boxrule=1.5pt,
    arc=4pt,
    breakable,
    enhanced
}

\newtcolorbox{warningbox}[1]{
    colback=white,
    colframe=accent,
    fonttitle=\bfseries\color{white},
    title=#1,
    boxrule=1.5pt,
    arc=4pt,
    breakable,
    enhanced
}

% Header and Footer
\pagestyle{fancy}
\fancyhf{}
\fancyhead[L]{\textcolor{secondary}{SIvision Academy -- C++ Course}}
\fancyhead[R]{\textcolor{secondary}{ESW Internship}}
\fancyfoot[C]{\thepage}
\renewcommand{\headrulewidth}{0.4pt}
\renewcommand{\footrulewidth}{0.4pt}

\begin{document}

% Cover Page
\begin{titlepage}
    \centering
    \vspace*{2cm}
    
    % Logo or Academic/Corporate Header
    {\scshape\Large SIvision Academy \par}
    \vspace{1.5cm}
    
    % Title
    {\Huge\bfseries C++ Programming Course \par}
    \vspace{0.5cm}
    {\large\bfseries ESW Internship Explanation Notes \par}
    \vspace{2cm}
    
    % Subtitle/Description
    \begin{tcolorbox}[colback=primary!5!white,colframe=primary,arc=8pt,halign=center]
        \centering\medskip
        A comprehensive compilation of code explanations, conceptual deep dives, and internship training material.
        \medskip
    \end{tcolorbox}
    
    \vfill
    
    % Author info
    {\large\bfseries Author:}\par
    {\large ESW Intern \par}
    
    \vspace{1.5cm}
    {\large\today \par}
\end{titlepage}

\newpage
\tableofcontents
\newpage

\section{Session 1: C++ Basics \& Memory Model}

\subsection{Pointer Casting and Memory Access Safety}

In C++, memory safety is paramount. Understanding how pointers interact with memory allocations on the stack is a crucial foundation of embedded systems software (ESW).

\subsubsection{Code Snippet}

\begin{lstlisting}[language=C++]
#include <iostream>
#include <stdint.h>
#include <string.h>
using namespace std;

void func() {
}

int main()
{
    uint8_t varX = 10;
    uint32_t* ptr = (uint32_t*)&varX;
}
\end{lstlisting}

\subsubsection{Detailed Explanation}

\begin{explanationbox}{Question 1: Effect of Using a 32-Bit Pointer to 8-Bit Data}
    A pointer's type determines how many bytes of memory are read or written when dereferencing it:
    \begin{itemize}
        \item \textbf{Dereference Size Mismatch}: The pointer \texttt{ptr} is of type \texttt{uint32\_t*}, meaning the CPU and compiler expect it to point to a \textbf{4-byte (32-bit)} memory block. However, \texttt{varX} is a \texttt{uint8\_t}, which only occupies \textbf{1 byte (8 bits)} of memory.
        \item \textbf{Out-of-Bounds Memory Access}: When dereferencing \texttt{*ptr} to read or write, the system reads or writes 4 full bytes starting from the address of \texttt{varX}. This accesses 3 adjacent bytes in memory that belong to other variables or stack frame metadata.
        \item \textbf{Stack Corruption (on Write)}: Attempting to assign a value (e.g., \texttt{*ptr = 0x12345678;}) will overwrite the 3 adjacent bytes on the stack, corrupting surrounding local variables or even the function return address, leading to a crash or security vulnerability.
        \item \textbf{Hardware Alignment Faults}: In embedded microcontrollers (like ARM Cortex-M), a 32-bit memory access must usually be aligned to a 4-byte boundary. Because \texttt{varX} is only 1 byte, it can be allocated at an unaligned address. Attempting a 32-bit access on an unaligned address triggers a hardware \textbf{UsageFault} or alignment abort.
        \item \textbf{Strict Aliasing Violations}: Compilers optimize code assuming that pointers of incompatible types do not alias (point to) the same memory. Violating this allows the compiler to generate unpredictable optimized machine code.
    \end{itemize}
\end{explanationbox}

\begin{explanationbox}{Question 2: Why is this Specific Casting Syntax Used?}
    In C++, the type system enforces strict type checking. Here is why the casting syntax \texttt{(uint32\_t*)\&varX} is necessary and what it does:
    \begin{itemize}
        \item \textbf{Address-of Operator (\&)}: \texttt{\&varX} retrieves the memory address where the 1-byte variable \texttt{varX} is stored. Its default type is \texttt{uint8\_t*} (pointer to a 1-byte unsigned int).
        \item \textbf{Type Mismatch Prohibition}: Without casting, assigning a \texttt{uint8\_t*} to a \texttt{uint32\_t*} directly fails to compile because they point to types of different sizes.
        \item \textbf{C-Style Casting}: The prefix \texttt{(uint32\_t*)} is a C-style cast. It acts as an override command to the compiler, saying: \textit{"Ignore the type mismatch and treat this address as if it points to a 32-bit integer."}
        \item \textbf{Modern C++ Alternative (\texttt{reinterpret\_cast})}:
        \begin{lstlisting}[language=C++]
uint32_t* ptr = reinterpret_cast<uint32_t*>(&varX);
        \end{lstlisting}
    \end{itemize}

    \textbf{Why C-Style Casts are Discouraged ("General and Easily Hidden"):}
    \begin{enumerate}
        \item \textbf{Lack of Intent Clarity (General)}: A C-style cast is a brute-force operator. When the compiler sees \texttt{(T)v}, it tries several casting mechanisms in sequence until one compiles:
        \begin{align*}
            \text{C-Style Cast } \longrightarrow \quad &\text{1. } \texttt{const\_cast} \\
            &\text{2. } \texttt{static\_cast} \\
            &\text{3. } \texttt{static\_cast} \text{ followed by } \texttt{const\_cast} \\
            &\text{4. } \texttt{reinterpret\_cast} \\
            &\text{5. } \texttt{reinterpret\_cast} \text{ followed by } \texttt{const\_cast}
        \end{align*}
        Because the compiler tries all of these, a developer cannot specify their exact intent. If type definitions change during code refactoring, a cast that was intended to be a safe \texttt{static\_cast} may silently turn into a dangerous \texttt{reinterpret\_cast} without the compiler raising any warnings.
        \item \textbf{Poor Searchability (Hidden)}: C-style casts are syntactically indistinguishable from regular parentheses. Searching a large codebase for instances of casting is extremely difficult because a simple text search matches function calls, mathematical expressions, and logic groupings. C++ cast keywords (like \texttt{reinterpret\_cast}, \texttt{static\_cast}) are unique, making them instantly searchable via tools like \texttt{grep}.
        \item \textbf{Static Analysis and Code Review Visibility}: Code quality tools (like Clang-Tidy or Cppcheck) can immediately flag the use of the \texttt{reinterpret\_cast} keyword as a high-risk operation. In code reviews, seeing the long, explicit \texttt{reinterpret\_cast} keyword immediately warns reviewers that type safety is being bypassed, prompting them to check for alignment, buffer sizes, and memory leaks. A C-style cast is easily overlooked.
    \end{enumerate}
\end{explanationbox}

\end{document}
